% Book 4, Year 2, chapter 11 -- The processor: datapath, control and pipelining.
% Provenance (never printed): sources/book4/BRIEFS.md, chapter 11.
\chapter{The Processor: Datapath, Control and Pipelining}\label{ch:b2:processor-and-pipelining}

A laundry has a washer, a dryer and a folding table, each taking half an hour. Done one at a time, four loads take
six hours. But as soon as the first load leaves the washer, the second can go in: with the three machines busy at
once, the four loads are done in three hours, and each extra load adds only half an hour. A processor does the same
with instructions. Here are five instructions going through a five-step processor, one row per instruction and one
column per clock cycle:

\omoutput[1-8]{code/bachelor-2/11-processor-and-pipelining/outputs/run_pipeline.txt}{Five instructions in 9
cycles instead of 25: F fetch, D decode, X execute, M memory, W write back.}\label{out:b2:processor-and-pipelining:hook}

No instruction finishes sooner than before; more of them finish per second. This chapter builds the processor that
runs the instructions of \cref{ch:b2:assembly-and-calling-convention}, first as one large circuit that does a whole
instruction in one clock cycle, then cut into five stages that work at once, and measures, on real programs run by
\code{rv}, what the overlap costs when instructions depend on one another.

\begin{recall}
Book 3: a multiplexer selects one of its inputs by a control value; a decoder turns a number into one active line;
an ALU computes an arithmetic or logic function chosen by its control inputs; every gate has a propagation delay.
A register holds its value until the clock edge; a clock paces a synchronous circuit; a finite-state machine changes
state on each edge. A processor is a datapath, which moves and transforms data, and a control unit, which drives it.
\cref{ch:b2:assembly-and-calling-convention}: the RV32I formats, the immediates, and branch offsets.
\end{recall}

\section{A single-cycle datapath}\label{sec:b2:processor-and-pipelining:datapath}

\begin{definition}[Single-cycle processor]\label{def:b2:processor-and-pipelining:datapath}
A \emph{single-cycle processor}\index{single-cycle processor} executes each instruction completely in one clock
cycle: in that cycle the instruction is fetched, its registers read, the ALU used, memory accessed and the result
written back, and the next clock edge stores the new program counter and the result.
\end{definition}

Its datapath, for a subset of RV32I (the register and immediate arithmetic, \code{lw}, \code{sw} and the six
conditional branches), needs these parts:
\begin{itemize}
\item the program counter (PC), a 32-bit register, and an instruction memory read at the PC;
\item the \omterm{def:b2:assembly-and-calling-convention:register}{register file}, with two read ports (\code{rs1}, \code{rs2}) and one write port (\code{rd});
\item the immediate generator, which gathers the scattered immediate bits of the I, S and B formats and sign-extends
them;
\item the ALU, and a data memory read or written at the address the ALU computes;
\item two adders for the next PC, $\mathit{pc} + 4$ and $\mathit{pc} + \mathit{imm}$;
\item multiplexers wherever two sources compete for one input.
\end{itemize}

\begin{omfigure}
\begin{tikzpicture}[font=\scriptsize, x=0.93cm, y=0.93cm,
  b/.style={draw, fill=omDef!8, align=center, inner sep=2pt},
  mx/.style={draw, fill=omProp!15, rounded corners=3pt, minimum width=3.2mm, minimum height=10mm, inner sep=0pt},
  sig/.style={text=omThm, font=\tiny}, w/.style={-{Stealth[length=1.6mm]}}]
\node[b, minimum width=7mm, minimum height=11mm] (pc) at (0,0) {PC};
\node[b, minimum width=15mm, minimum height=20mm] (im) at (2.0,0) {instruction\\ memory};
\node[b, minimum width=17mm, minimum height=24mm] (rf) at (5.2,0) {register\\ file};
\node[b, ellipse, inner sep=1pt] (ig) at (5.2,-2.3) {imm gen};
\node[mx] (m1) at (7.1,-0.3) {};
\node[b, minimum width=8mm, minimum height=17mm, fill=omThm!10] (alu) at (8.4,0.2) {ALU};
\node[b, minimum width=15mm, minimum height=20mm] (dm) at (10.4,-0.2) {data\\ memory};
\node[mx] (m2) at (12.2,-0.2) {};
\node[b] (a4) at (2.0,1.9) {$+4$};
\node[b] (ab) at (8.4,2.8) {$+$};
\node[mx] (m3) at (10.6,2.6) {};
\node[sig, above=0pt of m1] {ALUSrc}; \node[sig, above=0pt of m2] {MemtoReg};
\node[sig] at (10.6,3.45) {PCSrc $=$ Branch $\wedge$ taken}; \node[sig, above=0pt of rf] {RegWrite};
\node[sig, above=0pt of dm] {MemRead, MemWrite}; \node[sig] at (8.3,-0.95) {ALU control};
\draw[w] (pc.east) -- (im.west);
\draw (im.east) -- (3.5,0); \draw (3.5,1.1) -- (3.5,-2.3);
\draw[w] (3.5,0.8) -- node[above, font=\tiny] {rs1} (4.35,0.8);
\draw[w] (3.5,0.2) -- node[above, font=\tiny] {rs2} (4.35,0.2);
\draw[w] (3.5,-0.4) -- node[above, font=\tiny] {rd} (4.35,-0.4);
\draw[w] (3.5,-2.3) -- (ig.west);
\draw[w] (rf.east |- 0,0.7) -- (alu.west |- 0,0.7);
\draw (rf.east |- 0,-0.1) -- (6.6,-0.1); \draw[w] (6.6,-0.1) -- (m1.west |- 0,-0.1);
\draw (ig.east) -- (6.5,-2.3); \draw[w] (6.5,-2.3) -- (6.5,-0.55) -- (m1.west |- 0,-0.55);
\draw[w] (m1.east) -- (alu.west |- m1.east);
\draw[w] (alu.east |- 0,0.2) -- (dm.west |- 0,0.2);
\draw (6.6,-0.1) -- (6.6,-1.3); \draw[w] (6.6,-1.3) -- (9.4,-1.3) -- (9.4,-0.8) -- (dm.west |- 0,-0.8);
\draw (9.1,0.2) -- (9.1,-1.7) -- (11.8,-1.7); \draw[w] (11.8,-1.7) -- (11.8,-0.45) -- (m2.west |- 0,-0.45);
\draw[w] (dm.east |- 0,0.05) -- (m2.west |- 0,0.05);
\draw (m2.east) -- (12.6,-0.2) -- (12.6,-3.2) -- (3.9,-3.2) -- (3.9,-0.95); \draw[w] (3.9,-0.95) -- (4.35,-0.95);
\node[font=\tiny, anchor=north] at (8,-3.2) {write-back: the value written to \texttt{rd}};
\draw (pc.north) -- (0,3.1); \draw[w] (0,1.9) -- (a4.west); \draw[w] (0,3.1) -- (ab.west |- 0,3.1);
\draw[w] (6.5,-0.55) -- (6.5,2.6) -- (ab.west |- 0,2.6);
\draw (a4.east) -- (10.1,1.9); \draw[w] (10.1,1.9) -- (10.1,2.3) -- (m3.west |- 0,2.3);
\draw[w] (ab.east) -- (m3.west |- ab.east);
\draw (m3.east) -- (12.0,2.6) -- (12.0,4.0) -- (-0.8,4.0) -- (-0.8,0); \draw[w] (-0.8,0) -- (pc.west);
\end{tikzpicture}
\par\smallskip{\small The single-cycle datapath of \code{datapath.py}. Data flows left to right; the write-back and
the next PC return to the left. The three multiplexers choose the ALU's second input (register or immediate), the
value written back (ALU result or loaded word) and the next PC ($\mathit{pc} + 4$ or the branch target). The \omterm{def:b2:processor-and-pipelining:control}{control
signals}, in colour, come from the control unit; a comparator beside the ALU (not drawn) decides \emph{taken}.}
\end{omfigure}

The model \code{datapath.py} computes, for one instruction word and the current state of an \code{rv} machine,
what this circuit does in one clock cycle, without changing the machine; the caller then lets \code{rv} execute the
instruction and compares the two.

\omcode{code/bachelor-2/11-processor-and-pipelining/listings/datapath.py}{89}{108}{One cycle of the
datapath}\label{lst:b2:processor-and-pipelining:cycle}

\omcode{code/bachelor-2/11-processor-and-pipelining/listings/datapath.py}{40}{50}{The immediate generator: the
layouts of \cref{ch:b2:assembly-and-calling-convention} read backwards}

Each line of \code{cycle} is a wire or a box of the figure. Tracing four instructions is a matter of following the
values: \code{add t2, t0, t1} reads two registers, the ALU adds, the multiplexer passes the ALU result to the write
port. \code{sw t2, 12(sp)} reads \code{sp} and \code{t2}, the ALU adds 12 to \code{sp} as an address, and the data
memory stores \code{t2} there; nothing is written to a register. \code{lw t3, 12(sp)} computes the same address,
reads memory, and writes the loaded word, through the second multiplexer, to \code{t3}. \code{beq t2, t3, same}
compares the two registers; when they are equal, the PC multiplexer takes $\mathit{pc} + \mathit{imm}$.

\section{Control}\label{sec:b2:processor-and-pipelining:control}

\begin{definition}[Control signal]\label{def:b2:processor-and-pipelining:control}
A \emph{control signal}\index{control signal} is a wire from the control unit that selects a multiplexer's input,
enables a write, or chooses an ALU operation. In a \omterm{def:b2:processor-and-pipelining:datapath}{single-cycle processor} the control unit is a combinational
circuit whose input is the instruction's \omterm{def:b2:assembly-and-calling-convention:format}{opcode} (with \code{funct3} and \code{funct7} for the ALU).
\end{definition}

\begin{method}[Deriving the control table]\label{met:b2:processor-and-pipelining:control-table}
For each class of instruction, follow the path its data must take through the datapath, then set each signal:
\begin{enumerate}
\item RegWrite is 1 when the instruction has a destination register, else 0 (stores, branches);
\item ALUSrc is 1 when the ALU's second operand is the immediate (I-type, loads, stores), 0 for a register;
\item MemRead and MemWrite are 1 only for loads and stores respectively;
\item MemtoReg is 1 when the written value comes from memory (loads); when RegWrite is 0 it does not matter;
\item Branch is 1 for the conditional branches; the PC multiplexer takes the target when Branch is 1 and the
comparison holds;
\item the ALU operation is ``add'' for addresses, ``compare'' for branches, and read from the funct fields otherwise.
\end{enumerate}
A signal that does not matter for an instruction may be set either way; setting it to 0 makes the table simpler,
never a signal that writes something.
\end{method}

\omcode{code/bachelor-2/11-processor-and-pipelining/listings/datapath.py}{23}{29}{The control table, one row per
\omterm{def:b2:assembly-and-calling-convention:format}{opcode}}\label{lst:b2:processor-and-pipelining:control}

\code{alu_control} then picks the operation from \code{funct3}, and \code{funct7} separates \code{add} from
\code{sub} and \code{srl} from \code{sra}; an \code{addi} has no \code{funct7}, its bit 30 being part of the
immediate, which is why \code{subi} does not exist. The model's check runs \code{four.S} and compares every
instruction it covers with \code{rv}:

\omoutput{code/bachelor-2/11-processor-and-pipelining/outputs/run_datapath.txt}{The \omterm{def:b2:processor-and-pipelining:control}{control signals} of each
instruction, and the agreement with \code{rv}.}\label{out:b2:processor-and-pipelining:datapath}

The chapter's test runs the same check on every program of the chapter: the datapath and \code{rv} agree on every
instruction of the subset. The \code{ret} that ends \code{four.S} is a \code{jalr}, which needs one more path,
$\mathit{pc} + 4$ into the \omterm{def:b2:assembly-and-calling-convention:register}{register file}: \cref{exo:b2:processor-and-pipelining:12} adds it.

The trouble with this design is its clock. The cycle must be long enough for the slowest instruction, the load,
which goes through instruction memory, \omterm{def:b2:assembly-and-calling-convention:register}{register file}, ALU, data memory and back to the \omterm{def:b2:assembly-and-calling-convention:register}{register file}. Suppose those
take 200, 100, 200, 200 and 100\,ps: the clock period is at least 800\,ps, and an \code{add}, which does not need data
memory, waits for nothing while its 600\,ps of work end.

\section{The five-stage pipeline}\label{sec:b2:processor-and-pipelining:pipeline}

\begin{definition}[Instruction pipeline, pipeline stage, pipeline register]\label{def:b2:processor-and-pipelining:pipeline}
An \emph{instruction pipeline}\index{instruction pipeline} splits the execution of an instruction into steps, the
\emph{pipeline stages}\index{pipeline stage}, each taking one clock cycle, so that several instructions are in
progress at once, one per stage. Between two stages a \emph{pipeline register}\index{pipeline register} holds what
the next stage needs: the values computed so far and the \omterm{def:b2:processor-and-pipelining:control}{control signals} of the instruction.
\end{definition}

The classic RISC pipeline has five stages, cut where the single-cycle datapath has its large blocks: IF (fetch from
instruction memory), ID (decode and read the registers), EX (the ALU), MEM (data memory) and WB (write the
register). Its \omterm{def:b2:processor-and-pipelining:pipeline}{pipeline registers} are named after the stages they separate: IF/ID, ID/EX, EX/MEM, MEM/WB. With the
delays above, the clock period becomes that of the slowest stage, 200\,ps.

\begin{definition}[Latency, throughput]\label{def:b2:processor-and-pipelining:throughput}
The \emph{latency}\index{latency (of an instruction)} of an instruction is the time from its fetch to its
completion. The \emph{throughput}\index{throughput} of a processor is the number of instructions it completes per
unit of time, or per cycle.
\end{definition}

The pipeline does not lower the \omterm{def:b2:processor-and-pipelining:throughput}{latency}: in five stages of 200\,ps an instruction takes 1000\,ps, more than the
800\,ps of the single cycle. It raises the \omterm{def:b2:processor-and-pipelining:throughput}{throughput}: in the steady state, one instruction completes every cycle.

\begin{proposition}[Ideal gain of a pipeline]\label{prop:b2:processor-and-pipelining:speedup}
On a $k$-stage pipeline with no \omterm{def:b2:processor-and-pipelining:hazard}{hazard}, $n$ instructions take $k + n - 1$ cycles. With a cycle $k$ times shorter than
the unpipelined one, the time of the unpipelined machine divided by that of the pipeline is
\[
\frac{nk}{k + n - 1},
\]
which tends to $k$ as $n$ grows.
\end{proposition}

\begin{proof}
The first instruction leaves the last stage at the end of cycle $k$. Each later one enters the pipeline one cycle
after the one before it and so leaves one cycle after it: the $n$-th leaves at the end of cycle $k + n - 1$. The
unpipelined machine takes $n$ cycles of length $k$ (in units of the short cycle), so $nk$. Dividing,
$nk/(k + n - 1) = k / (1 + (k - 1)/n) \to k$.
\end{proof}

The hook's five instructions took $5 + 5 - 1 = 9$ cycles. In practice the stages are never perfectly balanced and
the \omterm{def:b2:processor-and-pipelining:pipeline}{pipeline registers} add their own delay, so the gain stays below $k$ even before the next section's \omterm{def:b2:processor-and-pipelining:hazard}{hazards}.

The chapter's timing model \code{pipeline.py} does not simulate wires: it takes the stream of instructions that
\code{rv} executed (one \code{Step} per instruction, with the registers it read and wrote) and computes the cycle in
which each enters each stage, with these rules: an instruction enters ID when the one before it leaves ID, and enters
EX when its operands are ready.

\omcode{code/bachelor-2/11-processor-and-pipelining/listings/pipeline.py}{36}{58}{The timing model, first part: when
each instruction reaches each stage}\label{lst:b2:processor-and-pipelining:schedule}

\section{Hazards}\label{sec:b2:processor-and-pipelining:hazards}

\begin{definition}[Hazard: structural, data, control]\label{def:b2:processor-and-pipelining:hazard}
A \emph{hazard}\index{hazard (pipeline)} is a situation in which the next instruction cannot enter its next stage in
the next cycle. It is \emph{structural} when two instructions need the same hardware in the same cycle; a \emph{data
hazard} when an instruction needs a value that an earlier instruction has not yet written; a \emph{control hazard}
when the processor does not yet know which instruction comes next, after a branch or a jump.
\end{definition}

Separate instruction and data memories remove the structural \omterm{def:b2:processor-and-pipelining:hazard}{hazard} between IF and MEM, and the \omterm{def:b2:assembly-and-calling-convention:register}{register file} is
written in the first half of a cycle and read in the second, so WB and ID can share a cycle. Data \omterm{def:b2:processor-and-pipelining:hazard}{hazards} are
everywhere: \code{add t0, a1, a2} followed by \code{sub t1, t0, a2} reads \code{t0} in ID while \code{add} is still
in EX. Without any help, \code{sub} must wait in ID until \code{add} has written \code{t0} in WB, two cycles later.

\begin{definition}[Stall, bubble, load-use hazard]\label{def:b2:processor-and-pipelining:stall}
A \emph{stall}\index{stall} holds an instruction in its stage for a cycle; the stage after it receives a
\emph{bubble}\index{bubble (pipeline)}, an empty slot that does nothing. A \emph{load-use hazard}\index{load-use
hazard} is a data \omterm{def:b2:processor-and-pipelining:hazard}{hazard} in which an instruction uses the register that the instruction just before it loads from
memory.
\end{definition}

\begin{definition}[Forwarding]\label{def:b2:processor-and-pipelining:forwarding}
\emph{Forwarding}\index{forwarding} (or \emph{bypassing}\index{bypassing}) sends a result to the ALU's input
directly from a \omterm{def:b2:processor-and-pipelining:pipeline}{pipeline register}, EX/MEM or MEM/WB, as soon as it has been computed, instead of waiting for it to be
written to the \omterm{def:b2:assembly-and-calling-convention:register}{register file} and read again.
\end{definition}

The value \code{add} computes is in EX/MEM at the end of its EX cycle, just in time for \code{sub}'s EX in the next
one. The \omterm{def:b2:processor-and-pipelining:forwarding}{forwarding} unit compares the source registers of the instruction in EX with the destinations of the two
instructions ahead of it:

\omcode{code/bachelor-2/11-processor-and-pipelining/listings/forwarding.py}{11}{16}{The \omterm{def:b2:processor-and-pipelining:forwarding}{forwarding} unit for one
source register}\label{lst:b2:processor-and-pipelining:forward}

\omoutput[1-8]{code/bachelor-2/11-processor-and-pipelining/outputs/run_forwarding.txt}{Where each operand comes
from, in a stream where each instruction uses the last
ones.}\label{out:b2:processor-and-pipelining:forwarding}

\begin{omfigure}
\begin{tikzpicture}[font=\scriptsize, st/.style={draw, fill=omDef!8, minimum width=14mm, minimum height=9mm},
  pr/.style={draw, fill=black!15, minimum width=2.5mm, minimum height=24mm, inner sep=0pt},
  fw/.style={-{Stealth[length=1.8mm]}, thick, omThm}]
\node[st] (if) at (0,0) {IF}; \node[pr] (r1) at (1.15,0) {}; \node[st] (id) at (2.3,0) {ID};
\node[pr] (r2) at (3.45,0) {}; \node[st, fill=omThm!10, minimum height=14mm] (ex) at (5.0,0) {ALU};
\node[pr] (r3) at (6.55,0) {}; \node[st] (mem) at (7.7,0) {MEM}; \node[pr] (r4) at (8.85,0) {};
\node[st] (wb) at (10.0,0) {WB};
\node[draw, fill=omProp!15, rounded corners=2pt, minimum width=3mm, minimum height=7mm] (ma) at (4.0,0.45) {};
\node[draw, fill=omProp!15, rounded corners=2pt, minimum width=3mm, minimum height=7mm] (mb) at (4.0,-0.45) {};
\foreach \r/\n in {r1/{IF/ID}, r2/{ID/EX}, r3/{EX/MEM}, r4/{MEM/WB}} \node[font=\tiny, rotate=90] at (\r) {\n};
\draw[->] (if) -- (r1); \draw[->] (r1) -- (id); \draw[->] (id) -- (r2);
\draw[->] (r2.east |- ma) -- (ma); \draw[->] (r2.east |- mb) -- (mb);
\draw[->] (ma.east) -- (ex.west |- ma.east); \draw[->] (mb.east) -- (ex.west |- mb.east);
\draw[->] (ex) -- (r3); \draw[->] (r3) -- (mem); \draw[->] (mem) -- (r4); \draw[->] (r4) -- (wb);
\draw[fw] (r3.north) -- ++(0,0.4) -| ([xshift=-0.6mm]ma.north);
\draw[fw, dashed] (r4.north) -- ++(0,0.7) -| ([xshift=0.6mm]ma.north);
\draw[fw] (r3.south) -- ++(0,-0.4) -| ([xshift=-0.6mm]mb.south);
\draw[fw, dashed] (r4.south) -- ++(0,-0.7) -| ([xshift=0.6mm]mb.south);
\end{tikzpicture}
\par\smallskip{\small The \omterm{def:b2:processor-and-pipelining:forwarding}{forwarding} paths. Each ALU input gets a multiplexer with three sources: the register value
read in ID, the result in EX/MEM (solid: one instruction ahead), or the result in MEM/WB (dashed: two ahead). The \omterm{def:b2:processor-and-pipelining:forwarding}{forwarding} unit drives them with the rule of
\cref{lst:b2:processor-and-pipelining:forward}.}
\end{omfigure}

\omterm{def:b2:processor-and-pipelining:forwarding}{Forwarding} cannot help a load: its value comes out of data memory at the end of MEM, one cycle later than an ALU
result. The model encodes these delays in three numbers: an instruction may be in EX one cycle after an ALU producer,
two after a load, and three after any producer when there is no \omterm{def:b2:processor-and-pipelining:forwarding}{forwarding} (it reads in ID, during the producer's
WB).

\begin{proposition}[One stall per load-use]\label{prop:b2:processor-and-pipelining:load-use}
With \omterm{def:b2:processor-and-pipelining:forwarding}{forwarding} from EX/MEM and MEM/WB, an instruction that uses the result of a load immediately before it is
stalled exactly one cycle, and an instruction that uses it two or more instructions later is not stalled by it.
\end{proposition}

The proof is \cref{exo:b2:processor-and-pipelining:8}. Here are both effects on \code{load_use.S}:

\omoutput[21-31]{code/bachelor-2/11-processor-and-pipelining/outputs/run_pipeline.txt}{With \omterm{def:b2:processor-and-pipelining:forwarding}{forwarding}: two \omterm{def:b2:processor-and-pipelining:stall}{stalls}
(*), one per load-use, and the stores and adds forwarded. A dot is an instruction waiting in IF because the stage
ahead is held.}\label{out:b2:processor-and-pipelining:load-use}

\omoutput[9-20]{code/bachelor-2/11-processor-and-pipelining/outputs/run_pipeline.txt}{The same program without
\omterm{def:b2:processor-and-pipelining:forwarding}{forwarding}: every dependence costs up to two cycles, 23 cycles
in all.}\label{out:b2:processor-and-pipelining:no-forwarding}

\begin{omfigure}
\begin{tikzpicture}[font=\scriptsize, x=6.2mm, y=5mm]
\foreach \c in {1,...,15} \node[text=black!60] at (\c,0.9) {\c};
\foreach \r/\t in {0/auipc a0, 1/addi a0, 2/lw t0, 3/add t1, 4/lw t2, 5/add t3, 6/sw t3, 7/addi a0, 8/jalr}
  \node[anchor=east, font=\scriptsize\ttfamily] at (0.4,-\r) {\t};
\foreach \r/\c/\s in {0/1/F,0/2/D,0/3/X,0/4/M,0/5/W, 1/2/F,1/3/D,1/4/X,1/5/M,1/6/W, 2/3/F,2/4/D,2/5/X,2/6/M,2/7/W,
  3/4/F,3/5/D,3/7/X,3/8/M,3/9/W, 4/5/F,4/7/D,4/8/X,4/9/M,4/10/W, 5/7/F,5/8/D,5/10/X,5/11/M,5/12/W,
  6/8/F,6/10/D,6/11/X,6/12/M,6/13/W, 7/10/F,7/11/D,7/12/X,7/13/M,7/14/W, 8/11/F,8/12/D,8/13/X,8/14/M,8/15/W}
  \node[draw, fill=omDef!10, minimum width=5.6mm, minimum height=4.2mm, inner sep=0pt] at (\c,-\r) {\s};
\foreach \r/\c in {3/6, 5/9} \node[draw, fill=omThm!25, minimum width=5.6mm, minimum height=4.2mm, inner sep=0pt] at (\c,-\r) {*};
\foreach \r/\c in {4/6, 6/9} \node[text=black!50] at (\c,-\r) {.};
\draw[-{Stealth[length=1.6mm]}, thick, omThm] (6,-2.3) -- (7,-2.7);
\draw[-{Stealth[length=1.6mm]}, thick, omThm] (9,-4.3) -- (10,-4.7);
\draw[-{Stealth[length=1.6mm]}, thick, omThm] (10.25,-5.3) -- (10.85,-5.75);
\end{tikzpicture}
\par\smallskip{\small The diagram of \cref{out:b2:processor-and-pipelining:load-use} drawn from the model's output.
The arrows are the forwarded values: from MEM/WB to an EX one cycle later (the two loads), and from EX/MEM (the
\code{add} to the \code{sw}). Each \omterm{def:b2:processor-and-pipelining:stall}{stall} slides every later instruction one column to the right.}
\end{omfigure}

The second half of the model handles control \omterm{def:b2:processor-and-pipelining:hazard}{hazards}. A conditional branch is resolved in EX, so when the processor
guesses wrong about it, the two instructions fetched after it are thrown away and fetching restarts after its EX:
two cycles lost. A \code{jal} knows its target in ID (it is $\mathit{pc} + \mathit{imm}$), one cycle lost; a
\code{jalr} needs a register, so EX, two cycles.

\omcode{code/bachelor-2/11-processor-and-pipelining/listings/pipeline.py}{59}{69}{The timing model, second part:
fetching again after a branch or a jump}

\section{Branch prediction}\label{sec:b2:processor-and-pipelining:prediction}

\begin{definition}[Branch prediction, predictor, misprediction penalty]\label{def:b2:processor-and-pipelining:prediction}
\emph{Branch prediction}\index{branch prediction} lets the processor fetch past a conditional branch before it is
resolved, along the direction a \emph{branch predictor}\index{branch predictor} guesses. When the guess is wrong,
the instructions fetched since are discarded (\emph{flushed}); the cycles lost are the \emph{misprediction
penalty}\index{misprediction penalty}, here 2.
\end{definition}

The simplest predictor says ``not taken'' for every branch, and is right whenever the branch falls through. Loops
defeat it: the branch that closes a loop is taken at every turn but the last. A dynamic predictor remembers what
each branch did. A one-bit predictor predicts that a branch does what it did last time; it is wrong twice per run of
a loop, at the last turn and again at the first turn of the next run.

\begin{definition}[Saturating counter]\label{def:b2:processor-and-pipelining:counter}
A two-bit \emph{saturating counter}\index{saturating counter} is a number from 0 to 3, incremented when the branch is
taken and decremented when it is not, without going past 0 or 3; it predicts ``taken'' when it is 2 or 3.
\end{definition}

\begin{omfigure}
\begin{tikzpicture}[font=\scriptsize, s/.style={draw, circle, minimum size=12mm, align=center, fill=omDef!8},
  t/.style={-{Stealth[length=1.8mm]}, bend left=30}]
\node[s] (0) at (0,0) {0\\ \tiny not taken}; \node[s] (1) at (2.6,0) {1\\ \tiny not taken};
\node[s, fill=omThm!12] (2) at (5.2,0) {2\\ \tiny taken}; \node[s, fill=omThm!12] (3) at (7.8,0) {3\\ \tiny taken};
\draw[t] (0) to node[above] {T} (1); \draw[t] (1) to node[above] {T} (2); \draw[t] (2) to node[above] {T} (3);
\draw[t] (3) to node[below] {N} (2); \draw[t] (2) to node[below] {N} (1); \draw[t] (1) to node[below] {N} (0);
\draw[-{Stealth[length=1.8mm]}] (0) edge[loop left] node {N} (0);
\draw[-{Stealth[length=1.8mm]}] (3) edge[loop right] node {T} (3);
\end{tikzpicture}
\par\smallskip{\small The two-bit predictor of one branch as a finite-state machine: T when the branch is taken, N when
it is not. A single surprise in state 3 or 0 moves the state but not the prediction.}
\end{omfigure}

\omcode{code/bachelor-2/11-processor-and-pipelining/listings/predictor.py}{29}{42}{A two-bit counter per branch,
kept by the address of the branch}\label{lst:b2:processor-and-pipelining:two-bit}

\begin{proposition}[Two bits for a loop]\label{prop:b2:processor-and-pipelining:two-bit}
Let the branch that closes an inner loop of $m \ge 2$ turns be taken $m - 1$ times and then not taken, at each
execution of the inner loop. Once its counter has reached 3 during one execution, a two-bit counter mispredicts it
exactly once per later execution of the inner loop; a one-bit predictor mispredicts it twice per execution.
\end{proposition}

The proof is \cref{exo:b2:processor-and-pipelining:11}. On the branch traces of three programs, the three predictors
give:

\omoutput[1-5]{code/bachelor-2/11-processor-and-pipelining/outputs/run_branches.txt}{Mispredicted branches out of
all conditional branches executed.}\label{out:b2:processor-and-pipelining:branches}

On \code{loop_sum.S} the two-bit counter is wrong twice: once at the first turn, while it warms up, once at the exit.
\code{collatz.S} is harder: its branch on the parity of $n$ follows the data, and there a one-bit predictor does
worse than predicting ``not taken'' at all, because an odd $n$ is always followed by an even one: each taken outcome
comes alone and costs the one-bit predictor two mistakes, the miss and the miss after it, where ``not taken'' misses
once. Real
processors keep thousands of counters, indexed by the address and by the outcomes of the last branches, and predict
far better than these three; the principle is the same.

\section{The performance equation}\label{sec:b2:processor-and-pipelining:performance}

\begin{definition}[Cycles per instruction, performance equation]\label{def:b2:processor-and-pipelining:cpi}
The \emph{cycles per instruction}\index{cycles per instruction} (CPI) of a run is its number of cycles divided by its
number of instructions. The \emph{performance equation}\index{performance equation} gives the time of a run:
\[
\text{time} = \text{instructions} \times \text{CPI} \times \text{cycle time}.
\]
\end{definition}

Each factor belongs to someone. The number of instructions depends on the program, the compiler and the instruction
set; the CPI on the pipeline's handling of \omterm{def:b2:processor-and-pipelining:hazard}{hazards} and on how often the program triggers them; the cycle time on the
circuit. A \omterm{def:b2:processor-and-pipelining:datapath}{single-cycle processor} has a CPI of 1 and a long cycle; the pipeline has a short cycle and a CPI above 1,
which \omterm{def:b2:processor-and-pipelining:forwarding}{forwarding} and prediction bring closer to 1. In the model, for $n$ instructions,
\[
\text{cycles} = n + 4 + \text{stall cycles} + \text{cycles lost to branches and jumps},
\]
an identity the chapter's test checks on every program and every configuration.

\omoutput[6-10]{code/bachelor-2/11-processor-and-pipelining/outputs/run_branches.txt}{The CPI of three programs in
three configurations of the pipeline.}\label{out:b2:processor-and-pipelining:cpi}

\begin{omfigure}
\begin{tikzpicture}
\begin{axis}[omrecord, ybar, width=0.9\linewidth, height=5.4cm, bar width=10pt, ymin=0, ymax=2.9,
    symbolic x coords={loop-sum, bubble-sort, collatz}, xtick=data, enlarge x limits=0.25,
    xticklabels={\texttt{loop\_sum}, \texttt{bubble\_sort}, \texttt{collatz}},
    ylabel={CPI}, legend style={font=\scriptsize, at={(0.5,1.03)}, anchor=south, legend columns=3},
    nodes near coords, every node near coord/.append style={font=\tiny, rotate=90, anchor=west}]
  \addplot[fill=black!25] table[x=program, y=nofwd] {code/bachelor-2/11-processor-and-pipelining/figures/cpi.dat};
  \addlegendentry{no forwarding}
  \addplot[fill=omDef!40] table[x=program, y=fwd] {code/bachelor-2/11-processor-and-pipelining/figures/cpi.dat};
  \addlegendentry{forwarding}
  \addplot[fill=omThm!50] table[x=program, y=fwdpred] {code/bachelor-2/11-processor-and-pipelining/figures/cpi.dat};
  \addlegendentry{forwarding + two-bit}
\end{axis}
\end{tikzpicture}
\par\smallskip{\small The CPI of \cref{out:b2:processor-and-pipelining:cpi}. \omterm{def:b2:processor-and-pipelining:forwarding}{Forwarding} helps every program; the
predictor helps the two loops, and hardly \code{collatz.S}, whose branches follow the data.}
\end{omfigure}

\omimg{book4/mips-r3000a-die.jpg}{0.5\linewidth}{The die of a MIPS R3000A, a 32-bit RISC processor, photographed on its
wafer. The squares along the edges are the pads wired to the package; the regular blocks are arrays of identical
cells, and the irregular regions are logic. Photograph: Pauli Rautakorpi, CC BY 3.0.}

\begin{inpractice}[Pipelines today]
Processors in desktops and phones have much longer pipelines than five stages, issue several instructions per cycle,
and execute them out of order when their operands are ready; the \omterm{def:b2:processor-and-pipelining:hazard}{hazards} of this chapter remain, handled by more
\omterm{def:b2:processor-and-pipelining:forwarding}{forwarding} paths, deeper buffers and better predictors. The \omterm{def:b2:processor-and-pipelining:cpi}{performance equation} still decides: a longer pipeline
shortens the cycle but raises the \omterm{def:b2:processor-and-pipelining:prediction}{misprediction penalty}. When code is slow on a modern machine, the usual suspects
are mispredicted branches on unpredictable data and loads that miss in the caches, the subject of
\cref{ch:b2:caches}.
\end{inpractice}

\section{Exercises}\label{sec:b2:processor-and-pipelining:exercises}

\begin{exercise}[$\star$]\label{exo:b2:processor-and-pipelining:1}
How many cycles does the five-stage pipeline take for 8 independent instructions? For $n$? What fraction of the
stage slots are used in the first case?
\end{exercise}

\begin{exercise}[$\star$]\label{exo:b2:processor-and-pipelining:2}
A designer copies the row of \code{lw} for \code{sw} and only clears MemRead and MemtoReg, leaving RegWrite at 1.
What does \code{sw t2, 12(sp)} then do besides storing? (Look at what bits 11 to 7 of an S-type word hold.)
\end{exercise}

\begin{exercise}[$\star$]\label{exo:b2:processor-and-pipelining:3}
Give the \omterm{def:b2:processor-and-pipelining:control}{control signals} for \code{lui} and the change it needs in the datapath. Add it to the model and check it
against \code{rv}.
\end{exercise}

\begin{exercise}[$\star$]\label{exo:b2:processor-and-pipelining:4}
Compute the ratio of \cref{prop:b2:processor-and-pipelining:speedup} for $k = 5$ and $n = 10$, $100$ and
$1000$, and for $k = 10$, $n = 10$. How many instructions are needed for a five-stage pipeline to reach 90\,\% of
its limit?
\end{exercise}

\begin{exercise}[$\star\star$]\label{exo:b2:processor-and-pipelining:5}
Draw the timing diagram, with \omterm{def:b2:processor-and-pipelining:forwarding}{forwarding}, of \code{exo_diagram.S}, and give its number of cycles.
\omcode{code/bachelor-2/11-processor-and-pipelining/listings/exo_diagram.S}{2}{6}{One load, used at once}
\end{exercise}

\begin{exercise}[$\star\star$]\label{exo:b2:processor-and-pipelining:6}
A \omterm{def:b2:processor-and-pipelining:forwarding}{forwarding} unit tests only \code{reg_write and rd == rs}, without \code{rd != 0}. Find a short program whose result
is then wrong, and say what it computes.
\end{exercise}

\begin{exercise}[$\star\star$]\label{exo:b2:processor-and-pipelining:7}
Add to the model a static predictor, ``backward taken, forward not taken'': a branch whose target lies before it is
predicted taken. Measure it on the three programs against the two-bit counter.
\end{exercise}

\begin{exercise}[$\star\star$]\label{exo:b2:processor-and-pipelining:8}
Prove \cref{prop:b2:processor-and-pipelining:load-use} from the model's rules.
\end{exercise}

\begin{exercise}[$\star\star$]\label{exo:b2:processor-and-pipelining:9}
Reorder the loop of \code{loop_sum.S} so that no instruction \omterm{def:b2:processor-and-pipelining:stall}{stalls} with \omterm{def:b2:processor-and-pipelining:forwarding}{forwarding}, without changing what it
computes, and measure the gain in cycles in the model.
\end{exercise}

\begin{exercise}[$\star\star\star$]\label{exo:b2:processor-and-pipelining:10}
In a program, 20\,\% of the instructions are conditional branches, mispredicted 10\,\% of the time with a penalty of 2
cycles; 25\,\% are loads, 40\,\% of which are followed at once by a user of the loaded value; there are no other
\omterm{def:b2:processor-and-pipelining:hazard}{hazards}. Compute the CPI. Which improvement is worth more: halving the misprediction rate, or reordering so that half
of the load-uses disappear?
\end{exercise}

\begin{exercise}[$\star\star\star$]\label{exo:b2:processor-and-pipelining:11}
Prove \cref{prop:b2:processor-and-pipelining:two-bit}. Does it hold for $m = 1$?
\end{exercise}

\begin{exercise}[$\star\star\star$]\label{exo:b2:processor-and-pipelining:12}
Add \code{jal} and \code{jalr} to the single-cycle datapath: the new paths, the new \omterm{def:b2:processor-and-pipelining:control}{control signals}, and the model,
checked against \code{rv} on a program that calls and returns.
\end{exercise}

\section{Problem: Making the sort run faster}\label{sec:b2:processor-and-pipelining:problem}

\begin{problem}[{Weekend problem --- Making the sort run faster}]\label{pb:b2:processor-and-pipelining:1}
A \omterm{def:b2:processor-and-pipelining:stall}{bubble} sort of 16 page counts is run through the pipeline model in its three configurations. The problem counts its
instructions, then its cycles, and reads the gain of each improvement in the \omterm{def:b2:processor-and-pipelining:cpi}{performance equation}.

\needspace{4\baselineskip}\medskip\noindent\textbf{Part I --- Counting instructions.}\par\nopagebreak
\omcode{code/bachelor-2/11-processor-and-pipelining/listings/bubble_sort.S}{25}{40}{The sort, a leaf}
\begin{enumerate}
\item How many instructions does one turn of the inner loop execute, without and with a swap?
\item How many does one turn of the outer loop execute outside the inner loop?
\item Prove that \code{sort} executes $\frac{5}{2}n(n-1) + 6(n-1) + 2S + 3$ instructions on $n \ge 1$ values, where
$S$ is the number of swaps.
\item Prove that $S$ is the number of inversions of the array, the pairs $i < j$ with $a_i > a_j$.
\item The array has 65 inversions and \code{main} executes 141 instructions. Check the total against the model's
count of 964.
\item How many loads and stores does the run execute?
\end{enumerate}

\needspace{4\baselineskip}\medskip\noindent\textbf{Part II --- Without forwarding.}\par\nopagebreak
\omoutput{code/bachelor-2/11-processor-and-pipelining/outputs/bubble_sort.txt}{The program's output.}
\begin{enumerate}[resume, beginpenalty=10000]
\item Which instructions of the inner loop depend on the instruction just before them?
\item Show that without \omterm{def:b2:processor-and-pipelining:forwarding}{forwarding} an instruction one place after its producer waits two cycles, and one two places
after waits one.
\item The model counts 552 \omterm{def:b2:processor-and-pipelining:stall}{stall} cycles. Account for them: 120 turns of the inner loop, 16 of the printing loop, 15
outer turns, and 5 single \omterm{def:b2:processor-and-pipelining:stall}{stalls} of two cycles each.
\item Under ``not taken'', which branches cost cycles, and how many cycles do they and the jumps lose in all?
\item Compute the cycles and the CPI, and check them against \cref{out:b2:processor-and-pipelining:cpi}.
\item Which costs more here, the data \omterm{def:b2:processor-and-pipelining:hazard}{hazards} or the control \omterm{def:b2:processor-and-pipelining:hazard}{hazards}?
\end{enumerate}

\needspace{4\baselineskip}\medskip\noindent\textbf{Part III --- Forwarding and prediction.}\par\nopagebreak
\begin{enumerate}[resume, beginpenalty=10000]
\item With \omterm{def:b2:processor-and-pipelining:forwarding}{forwarding}, which dependences still \omterm{def:b2:processor-and-pipelining:stall}{stall}, and how many cycles do they cost?
\item Propose a change of the inner loop that removes them, and say what else it saves.
\item The two-bit counters mispredict 60 branches: 41 of \code{bge t4, t3, keep}, 16 of the inner loop's
\code{bltu}, 2 of the printing loop's, 1 of the outer test. Explain each number.
\item Why is the swap test the hardest to predict?
\item How many cycles do branches and jumps now lose?
\item Compute the cycles and the CPI.
\end{enumerate}

\needspace{4\baselineskip}\medskip\noindent\textbf{Part IV --- The performance equation.}\par\nopagebreak
\begin{enumerate}[resume, beginpenalty=10000]
\item With the stage delays of this chapter, the single-cycle clock is 800\,ps and the pipeline's 200\,ps. Compute the
time of the run on the \omterm{def:b2:processor-and-pipelining:datapath}{single-cycle processor} and on the pipeline in its three configurations.
\item How many times faster than the \omterm{def:b2:processor-and-pipelining:datapath}{single-cycle processor} is the best pipeline? Why not 4?
\item Which factor of the \omterm{def:b2:processor-and-pipelining:cpi}{performance equation} does each of the three improvements (pipelining, \omterm{def:b2:processor-and-pipelining:forwarding}{forwarding},
prediction) act on?
\item What would the CPI be with a perfect predictor?
\item And with a perfect predictor and the change of question 14?
\item State the result of the problem.
\end{enumerate}
\end{problem}
